\section{实验方法}  % 原来是 \chapter

\subsection{网络结构设计}

本实验构建了一个三层全连接神经网络。

网络结构：第一层全连接层将输入从 785 维映射到 256 维，后接 ReLU 激活函数；第二层全连接层将 256 维映射到 128 维，后接 ReLU 激活函数；第三层全连接层将 128 维映射到 10 维，后接 Softmax 函数输出各类别概率。

权重初始化：采用 He 初始化方法，各层权重按 $\mathrm{scale} = \sqrt{2.0 / \mathrm{input\_dim}}$ 设置随机初始值。

模块实现：将矩阵乘法、ReLU 激活、Softmax 输出、交叉熵损失分别封装为独立类，每个类均实现 forward 前向传播和 backward 反向传播方法，便于梯度计算与参数更新。

\subsection{网络结构示意图}

\begin{figure}[!htbp]
\centering
\begin{tikzpicture}[node distance=1.5cm]
\node (input) [draw, rectangle] {输入层 (784)};
\node (fc1) [draw, rectangle, below of=input] {全连接层 W1: 785×256};
\node (relu1) [draw, rectangle, below of=fc1] {ReLU};
\node (fc2) [draw, rectangle, below of=relu1] {全连接层 W2: 256×128};
\node (relu2) [draw, rectangle, below of=fc2] {ReLU};
\node (fc3) [draw, rectangle, below of=relu2] {全连接层 W3: 128×10};
\node (softmax) [draw, rectangle, below of=fc3] {Softmax};
\node (loss) [draw, rectangle, below of=softmax] {CrossEntropy Loss};

\draw[->] (input) -- (fc1);
\draw[->] (fc1) -- (relu1);
\draw[->] (relu1) -- (fc2);
\draw[->] (fc2) -- (relu2);
\draw[->] (relu2) -- (fc3);
\draw[->] (fc3) -- (softmax);
\draw[->] (softmax) -- (loss);
\end{tikzpicture}
\caption{网络结构示意图}
\label{fig:network}
\end{figure}

\subsection{各模块说明}

\subsubsection{Matmul（矩阵乘法层）}

\textbf{前向传播}：
\[
y = xW
\]

\textbf{反向传播}：
\[
\frac{\partial L}{\partial x} = \frac{\partial L}{\partial y} W^T
\]
\[
\frac{\partial L}{\partial W} = x^T \frac{\partial L}{\partial y}
\]

\subsubsection{ReLU 激活函数}

\[
f(x) = \max(0, x)
\]

\textbf{梯度}：
\[
f'(x) = 
\begin{cases} 
1 & x > 0 \\
0 & x \leq 0
\end{cases}
\]

\subsubsection{Softmax}

\[
S_i = \frac{e^{x_i}}{\sum_j e^{x_j}}
\]

\subsubsection{交叉熵损失函数}

\[
L = -\sum y \log(\hat{y})
\]

\subsection{优化方法}

\subsubsection{L2 正则化}

\[
L = L + \lambda \|W\|^2
\]

\subsubsection{动量梯度下降}

\[
v = \mu v + \eta \nabla W
\]
\[
W = W - v
\]

\subsubsection{学习率衰减}

每 10 个 epoch：
\[
lr = lr \times 0.95
\]
